\textbf{How Cellpose works.} Cellpose \citep{stringer2021} was designed to
separate touching cells in microscope images, which is the same problem as
separating packed froth bubbles. Instead of drawing boxes, it predicts for
every pixel which way to move to reach the centre of the object it belongs
to (Figure~\ref{fig:arch-froth}).

\begin{figure}[htbp]
\centering
\begin{tikzpicture}[
  font=\footnotesize, >=Stealth, node distance=4mm,
  blk/.style={draw, rounded corners=1.5pt, align=center, fill=gray!8, minimum height=10mm, inner sep=3pt},
  io/.style={blk, fill=white, dashed},
  arr/.style={->, thick}]
\node[io] (in) {Froth image\\(in-focus band)\\tiles of\\256$\times$256};
\node[blk, right=of in] (pe) {Patch\\embedding\\8$\times$8 px\\patches};
\node[blk, right=of pe] (vit) {SAM image encoder\\ViT-L\\24 transformer blocks\\(global attention)};
\node[blk, right=of vit] (head) {Output layer\\$\rightarrow$ full-resolution\\maps};
\node[blk, below=9mm of head] (maps) {3 maps per pixel:\\horizontal flow,\\vertical flow,\\object probability};
\node[blk, left=of maps] (dyn) {Flow tracking\\pixels follow the flows\\and gather at centres};
\node[io, left=of dyn] (out) {One mask per bubble\\$\rightarrow$ equivalent-area\\diameter $\rightarrow$ D32};
\draw[arr] (in) -- (pe); \draw[arr] (pe) -- (vit); \draw[arr] (vit) -- (head);
\draw[arr] (head) -- (maps); \draw[arr] (maps) -- (dyn); \draw[arr] (dyn) -- (out);
\end{tikzpicture}
\caption{Cellpose-SAM as used for froth bubbles, with default settings and no froth training.}
\label{fig:arch-froth}
\end{figure}

\begin{enumerate}
\tightlist
\item \textbf{Encoder.} Cellpose-SAM \citep{pachitariu2025} replaces the
original Cellpose U-Net with the image encoder of Segment Anything
\citep[SAM,][]{kirillov2023}, a vision transformer (ViT-L) pretrained on
about one billion object masks. The image is split into 8$\times$8~px
patches, and the transformer blocks let every patch compare itself with
every other patch in the tile. SAM's own prompt and mask decoders are
removed.
\item \textbf{Outputs.} For every pixel the network predicts a horizontal
and a vertical flow (a small arrow pointing towards the centre of the
object) and the probability that the pixel belongs to any object.
\item \textbf{Flow tracking.} Pixels with high probability are moved step
by step along the arrows. Pixels that end up at the same point form one
object, so two bubbles that touch are still split along the wall between
them. Masks whose flows are inconsistent are discarded.
\end{enumerate}

The released model was trained on about 23,000 microscope images with 3.3
million outlined cells \citep{pachitariu2025}. It was used here exactly as
released, which is why froth measurement needed no labelled froth training
data.

